\section{From amplitudes to cut integral families}
\label{sec:cut-families}

\subsection{Ordered amplitude interferences}

At a fixed perturbative order, the elementary object is an ordered pair
\begin{equation}
 \mathcal W_{ij}
 =\sum_{\text{colors, spins}}
 \mathcal M_i\,\mathcal M_j^*.
 \label{eq:ordered-interference}
\end{equation}
The ordering matters whenever the two sides contain virtual loops, since the
ordinary propagators of $\mathcal M_i$ and $\mathcal M_j^*$ carry opposite
causal boundary values.  Ordinary spin and polarization sums are performed
first.  Spin sums on collinear parton lines associated with an incoming or
identified hadron are then replaced by the correlators in
Eqs.~\eqref{eq:pdf-correlator} and \eqref{eq:ff-correlator}.  This gives a
Dirac- and color-contracted numerator in $D=4-2\eps$ dimensions while
preserving the declared polarization channel.

For $m$ unobserved real partons with momenta $k_r$, the physical phase
space is
\begin{equation}
 \dd\Phi_m(q;\{k_r\})
 =(2\pi)^D\delta^{(D)}\!\left(q-\sum_{r=1}^m k_r\right)
 \prod_{r=1}^m
 \frac{\dd^Dk_r}{(2\pi)^{D-1}}\delta_+(k_r^2),
 \label{eq:physical-m-particle-phase-space}
\end{equation}
where $\delta_+(k^2)=\theta(k^0)\delta(k^2)$.  One phase-space momentum is
eliminated with the momentum-conservation delta function before families are
constructed.  The remaining positive-energy constraints are retained as
distinguished cut lines.

\subsection{Reverse unitarity and causal sectors}

For a cut momentum $q_c$, denominator $D_c=q_c^2-m_c^2$, and energy
orientation $\xi_c=\pm1$, define
\begin{equation}
 \mathcal C_{\xi_c}^{(n)}(D_c)
 =\frac{\theta(\xi_cq_c^0)}{2\pi i}
 \left[\frac{1}{(D_c-i0)^n}-\frac{1}{(D_c+i0)^n}\right]
 =\theta(\xi_cq_c^0)
 \frac{(-1)^{n-1}}{(n-1)!}\delta^{(n-1)}(D_c).
 \label{eq:oriented-cut}
\end{equation}
Reverse unitarity treats these distributions as propagator slots under IBP
differentiation~\cite{Anastasiou:2002yz}.  In particular, a squared cut is a
derivative of the same on-shell constraint; it is not a second cut particle.
The cut-family condition is
\begin{equation}
 I(\nu_1,\ldots,\nu_N)=0
 \quad\text{if}\quad
 \nu_c\leq0\ \text{for any designated cut slot }c.
 \label{eq:cut-family-condition}
\end{equation}

Three causal cases are kept distinct.  A virtual loop belonging to the
amplitude carries the Feynman prescription, while a virtual loop belonging to
the conjugate amplitude carries its complex conjugate.  A denominator that
contains virtual momenta from both sides is not assigned a prescription by
algebraic guesswork.  The present construction rejects such a family unless a
separate contour derivation is supplied.  An ordinary denominator containing
only external and real phase-space momenta needs no $i0$ label when its sign
is fixed on the closed physical phase space.  Appendix~\ref{app:fixed-sign}
gives a finite criterion for massless two-beam kinematics.  A zero on a soft
or collinear boundary is allowed; denominators whose sign changes in the
interior are rejected.

\subsection{Cut-preserving partial fractions}

The literal propagators of one ordered interference need not be linearly
independent as polynomials in the scalar products of the integration
momenta.  Suppose
\begin{equation}
 c_0+\sum_{j=1}^{N}c_jD_j=0,
 \label{eq:denominator-linear-relation}
\end{equation}
with coefficients independent of the integration momenta.  Standard
multivariate partial fractions decompose the product of denominators into
terms containing independent subsets~\cite{Feng:2012iq}.  For cut integrals,
the algebra is performed subject to Eq.~\eqref{eq:cut-family-condition}: every
surviving term must contain every physical cut with positive power.  This is
why cut propagators cannot be treated as ordinary factors and filtered only
after partial fractioning.

The order of the construction is
\begin{equation}
 \text{literal denominators}
 \longrightarrow
 \text{cut-preserving partial fractions}
 \longrightarrow
 \text{independent denominator subsets}
 \longrightarrow
 \text{complete scalar-product families}.
 \label{eq:family-construction-order}
\end{equation}
Only in the last step are irreducible scalar products added as auxiliary
slots.  This avoids choosing a process-specific expected family before the
amplitude algebra has shown which denominators actually occur.  Any new
auxiliary denominator that contains no virtual loop momentum is subjected to
the same fixed-sign criterion as a literal real-emission denominator.

\subsection{BMHV tensors and dimension recurrences}

Polarized traces are evaluated in the BMHV scheme
\cite{Breitenlohner:1977hr}.  The metric and every integration momentum split
as
\begin{equation}
 g_D^{\mu\nu}=\bar g^{\mu\nu}+\widehat g^{\mu\nu},
 \qquad
 \ell_r^\mu=\bar\ell_r^\mu+\widehat\ell_r^\mu,
 \qquad
 \bar g^\mu{}_{\mu}=4,
 \qquad
 \widehat g^\mu{}_{\mu}=D-4.
 \label{eq:bmhv-split-main}
\end{equation}
External hadron momenta, light-cone directions, and spin vectors are declared
four dimensional.  Evanescent scalar products therefore occur only among
integration momenta.  Tensor integration converts a monomial of total
evanescent rank $2r$ into scalar integrals in $D+2r$ dimensions.  Tarasov
dimension recurrences then express those shifted-dimensional integrals in the
same $D$-dimensional families~\cite{Tarasov:1996br}.  The recurrence is
performed before the final IBP reduction, so the objects submitted to
\Kira\ are ordinary $D$-dimensional scalar integrals with the original cut
slots intact.  Appendix~\ref{app:bmhv-dimension-shift} states the normalization
and the checks needed for this step.
